\chapter{物理仿真与动力学}

\paragraph*{本章阅读锚点}
物理方向的 3DGS 不能只按动态新视角合成来读。原始 3DGS 只保证多视图图像一致，得到的高斯可能视觉上正确，却不一定有质量、材料、碰撞、流体边界或可仿真的拓扑。物理方法的核心问题是：如何把高斯从“可渲染外观原语”变成“可交互、可模拟、可预测的物质表示”。本章的技术路线大致分为四类：第一类把物理属性或物理目标接到 3DGS 优化中；第二类用网格、弹簧质点或杆模型补上可变形物体的结构；第三类把流体、火焰等粒子或体积动力学纳入高斯场；第四类针对风驱动物体等具体物理现象做可微建模。

\section{物理仿真集成}

这一节的锚点是“视觉重建不等于物理可用”。标准 3DGS 优化的是渲染损失，隐含地允许空洞、悬浮、高斯漂浮、内部封闭、材料错误等问题；这些问题在普通新视角指标中可能不明显，但一旦要倒水、飞行、抓取、碰撞或变形，就会失败。PUGS 从视觉和基础模型中估计物理属性，PG-3DGS 把可微物理目标直接加入 3DGS 优化，GASP/PIDG/PhysFlow/Pixie 则分别把黑盒物理引擎、本构方程、视频扩散和监督物理预测接入高斯表示。



% METHOD: 2409.05819
\paragraph{GASP（把高斯交给黑盒物理引擎）（2024.09）}
GASP~\cite{2409_05819} 的锚点是“先把 3DGS 转成网格，再用传统物理引擎”。这种路线依赖高质量 meshing，而从任意高斯场提取平滑、闭合、可仿真的网格本身并不稳定。GASP 的改变是采用 GaMeS/flat GS 一类扁平高斯表示，把每个高斯参数化成可导出伪网格或三角片的形式，使物理引擎看到的是普通三角面、顶点或 pseudo-mesh，而渲染端仍保留高斯表示。

具体流程先用 GaMeS 训练 \texttt{gs\_flat} 模型，再运行伪网格生成脚本，把扁平高斯转成仿真文件：Taichi 路线使用顶点张量，Blender 和 Genesis 路线使用可导入的 \texttt{obj}。如果场景中高斯数量过大，还可以先生成层次结构，只让核心高斯或簇代表进入仿真，再把仿真位移映射回子高斯。物理引擎完成碰撞、破碎、弹性或 MPM 更新后，GASP 再用更新后的伪网格状态驱动高斯渲染。相对 PG-3DGS 的可微物理联合优化，GASP 更实用地强调“与已有黑盒物理引擎集成”：它不要求重写模拟器，也不要求物理梯度端到端回传，但需要 flat GS/pseudomesh 这一中间层维持高斯与仿真状态的一致。
\textbf{实现细节：} 与紧耦合路线将高斯基元直接作为 MPM 粒子不同，GASP（2024.09） 修改了高斯基元与物理引擎的对接方式——用代理点结构替代高斯基元直接参与仿真，从而将物理仿真从 MPM 单一引擎解放为任意黑盒物理引擎。紧耦合路线的局限在于：PhysGaussian 等方法需要修改 MPM 内核以适应高斯参数（如用一阶 Taylor 近似传播变形梯度），因此无法接入 Blender、Genesis 等通用仿真软件。GASP 的核心思路是：既然物理引擎只关心点的运动，那就只把点的位置交给物理引擎，再把运动结果映射回高斯参数——引擎是黑盒，GASP 不需要知道引擎内部如何计算。
\textbf{实现补充：} 实现上还要核对输入、输出、参数、特征、初始化、训练损失、正则、渲染或推理接口，以及优化中的剪枝、分裂、采样、编码或解码步骤；这些抓手决定该方法到底改了 3DGS 管线的哪一环。
\textbf{方法信息：} 代码：\url{https://github.com/waczjoan/GASP}，约 60 stars。
% METHOD: 2411.14423
\paragraph{PhysFlow（基础模型与视频扩散引导的物理仿真）（2024.11）}
PhysFlow~\cite{2411_14423} 的锚点是 PhysGaussian/PhysDreamer 这类从视觉场景生成物理动态的方法。它的问题判断是：单靠输入图像很难知道材料类型和参数，单靠渲染损失或 SDS 又容易显存高、噪声大、难以捕捉大幅运动。PhysFlow 的改变是把多模态基础模型、3DGS 场景表示、可微 MPM 和视频扩散生成的运动参考结合起来。

具体流程先根据输入形式选择不同重建器：多视角可用 PGSR，动态视频可用 Grid4D，单图可用 Splatt3R，得到一个 3DGS 场景。然后用 GPT-4 等多模态模型识别材料类型、密度和初始参数，覆盖弹性、橡皮泥、金属、泡沫、沙土、牛顿流体和非牛顿流体等类别。物理演化由可微 MPM 执行，高斯场和材料属性共同进入模拟。为了优化材料参数，PhysFlow 不直接用视频扩散图像做 noisy render loss，而是从扩散生成的视频中提取光流，把运动趋势作为监督信号；同时用 3DGS 渲染中间帧帮助阈值和几何对齐。它相对 PG-3DGS 的差异是更依赖基础模型和生成式先验来补材料与运动参考，适合输入信息不足但希望自动生成物理动态的场景。
\textbf{实现细节：} PhysFlow（2024.11） 修改了材料参数优化的监督信号环节——用光流损失（optical flow loss）替代传统的渲染损失（render loss）和分数蒸馏采样损失（SDS loss），实现从任意输入（单图/多视角/视频）自动估计材料属性并进行 4D 物理仿真。Baseline 做法与问题。在 PhysFlow 之前，基于视频扩散的材料参数优化主要有两条路线。另一条路线（如 Physics3D）采用 SDS 损失 $L_{\mathrm{sds}}$，从预训练扩散模型中蒸馏梯度，但 SDS 需要对整个扩散模型做反向传播，显存开销极大，且在结合复杂文本提示的视频扩散模型时更为严重。此外，上述方法均局限于弹性材料，无法处理流体、沙子等复杂材料。光流引导的参数优化。PhysFlow 的核心洞察是：优化材料参数并不需要逐像素对比图像，而只需对比运动场——光流恰好编码了像素级位移信息，是材料参数对运动的"直接传感器"。具体流程为：首先利用 GPT-4 从输入图像推断材料类型 $M$、密度 $\rho$ 和初始参数 $M_p^{}0$；然后根据输入类型（多视角/视频/单图）分别选择 PGSR、Grid4D 或 Splatt3R 构建 3D 高斯；接着通过视频扩散模型（根据输入类型可选 SVD、CogVideoX、MotionCtrl 等）生成参考视频，并用 RAFT 计算相邻帧光流 $U(I_t, I_{t+1})$；最终在可微分 MPM（Material Point Method，一种混合拉格朗日-欧拉的粒子仿真方法，将材料离散为粒子追踪变形，在背景网格上求解动量守恒）框架下，最小化光流损失： 其中 $\hat{U}$ 是仿真渲染帧之间的光流。直观地说，该损失迫使仿真运动与参考视频中的运动对齐，而非逐像素对比外观——这避免了图像级噪声的干扰，同时使梯度信号更直接地反映材料参数对运动的影响。材料类型覆盖。
\textbf{方法信息：} 代码：\url{https://github.com/zhuomanliu/PhysFlow}；引用很多，Semantic Scholar 约 35 次。



% METHOD: 2502.12231
\paragraph{PUGS（零样本物理属性理解）（2025.02）}
PUGS~\cite{2502_12231} 的锚点是 NeRF2Physics 式的“从多视图外观估计物理属性”。它相对原始 3DGS 的问题判断很直接：高斯为了渲染可以漂浮或错位，但质量、摩擦、硬度这类物理量需要几何体积和局部材料分布稳定。PUGS 的改变是把 3DGS 重建、视觉语言模型的常识推理、以及基于高斯体积的物理量积分连成一条零样本物理理解管线。

具体流程分三步。第一步是 shape-aware 3DGS reconstruction：在普通 RGB 重建损失外加入几何感知正则，让高斯分布更贴合物体实际空间结构，减少为了渲染而产生的漂浮原语；同时用 SAM 分割和对比学习训练 region-aware features，使相同材料或部件的高斯在特征空间更接近。第二步是 VLM-based physical property prediction：把图像或局部区域送入视觉语言基础模型，利用其材料常识预测质量、摩擦、硬度等物理属性。第三步是 feature-based property propagation：把 VLM 给出的区域级属性沿着高斯的局部特征传播到所有高斯，再通过高斯体积积分得到物体级质量等结果。它相对“只问 VLM 物体有多重”的方法更可靠，因为几何体积来自 3DGS；相对标准 3DGS 更物理，因为每个高斯不再只有颜色和不透明度，也承载可传播的物理属性。
\textbf{实现细节：} 沿"监督学习/前馈"路线，PIXIE 是首个通过监督学习训练通用神经网络预测 3D 场景物理属性的前馈框架，利用 CLIP 预训练特征把外观与物理特性关联；ReconPhys 是首个从单目视频联合推断物理属性和 3DGS 的前馈框架，推理速度小于 1 秒。沿"概率/不确定性"路线，PhysGS 把物理属性估计建模为 3DGS 上的贝叶斯推理问题，建模认知不确定性和偶然不确定性。仿真效率方面，FastPhysGS 提出实例感知粒子填充和双向图解耦优化，在 1 分钟内完成高保真物理仿真；PIG 通过 2D 像素到 3D 高斯的快速精确映射实现多材料耦合交互。
\textbf{方法信息：} 代码：\url{https://github.com/evernorif/pugs}；GitHub 关注度较高，约 105 stars。



% METHOD: 2508.17437
\paragraph{Pixie（从像素监督学习 3D 物理）（2025.08）}
Pixie~\cite{2508_17437} 与 PUGS 同样关心物理属性，但技术路线从零样本推理转向监督学习。PUGS 主要依赖 VLM 常识和少量结构正则；Pixie 的改变是构建成对的 3D 资产与物理标注数据集 PixieVerse，并训练一个前馈网络，从 3D 视觉特征直接预测逐点材料属性。

具体做法先把 2D 预训练视觉特征提升到三维，例如用 NeRF 或高斯场蒸馏 CLIP 等特征，得到覆盖物体空间的 semantic feature field。随后，一个 U-Net 式网络从这些 3D 特征预测材料类别、本构定律、杨氏模量、泊松比和密度等物理参数。预测出的材料场再与 3DGS 场景表示结合，用 MPM 做下游物理仿真。它相对 PhysFlow 的优势是推理快，测试时不需要为每个场景长时间优化材料；相对 PUGS 的差异是把常识推理变成大规模监督学习，期望通过合成物理标注获得更稳定的泛化。
\textbf{实现细节：} PIXIE（2025.08） 更进一步，提出首个通过监督学习训练通用神经网络预测 3D 场景物理属性的前馈框架——将属性预测从"每次推理"提升为"一次训练、处处适用"。其核心洞察是利用 CLIP 等预训练视觉特征将外观与物理特性关联——金属看起来光滑且反光，橡胶看起来粗糙且哑光。构建了 PIXIEVERSE 数据集（最大的 3D 资产与物理标注配对数据集），比测试时优化方法快数个数量级且效果提升 1.46--4.39 倍。在 PIXIE 将属性预测从优化范式转向前馈范式之后，后续工作沿加速和输入简化两个方向展开。
\textbf{锚点细化：} PIXIE：针对 Sim Anything 依赖 MLLM 测试时优化、单场景不泛化的局限，用\textbf{监督学习训练的通用前馈神经网络}替代逐场景 MLLM 推理，核心思想是\textbf{从``每次推理''提升为``一次训练、处处适用''}——利用 CLIP 预训练视觉特征将外观与物理特性关联（金属反光、橡胶哑光），构建 PIXIEVERSE 数据集（最大 3D 资产与物理标注配对），比测试时优化快数个数量级且效果提升 1.46--4.39 倍。
\textbf{方法信息：} Adobe 参与了这篇论文；代码：\url{https://github.com/vlongle/pixie}；GitHub 关注度高，约 295 stars。



% METHOD: 2511.06299
\paragraph{PIDG（本构约束的可变形高斯）（2025.11）}
Physics-Informed Deformable Gaussian Splatting~\cite{2511_06299} 的锚点是动态 4DGS，但它不满足于学习一个黑盒变形场。普通 4DGS 可以把每个时间的外观拟合好，却不保证运动符合材料力学；PIDG 的改变是把每个高斯看作拉格朗日物质点，让它带有速度、应力和时变本构参数，并用 Cauchy 动量方程残差约束时空运动。

具体做法先用 4D 分解哈希编码表示时空场：空间网格负责静态结构，多个时空哈希网格负责时间变化，再用空间 MLP 的注意力权重调制时间编码，分离场景中的静态和动态部分。对每个高斯粒子，模型独立预测速度和本构应力，物理损失来自 Cauchy momentum residual，要求预测运动与材料场满足连续介质力学约束。为了让单目动态视频也能收敛，方法把三维粒子运动投影到二维，与光流监督匹配，形成拉格朗日粒子流约束。相对 PG-3DGS 面向具体物理任务，PIDG 更像给动态高斯加上统一材料规律，让重建出的运动不仅可看，而且在局部力学上更一致。
\textbf{实现细节：} PIDG（2025.11） （Physics-Informed Deformable Gaussian Splatting）修改了动态 3DGS 的变形场模块——用物理定律约束的连续介质力学变形替代纯数据驱动的 MLP 变形——渲染管线和哈希编码框架保持不变。针对纯 MLP 变形场忽略材料本构律、仅靠光度损失无法约束三维粒子运动的问题，PIDG 将每个高斯建模为携带速度与应力张量的拉格朗日粒子，通过 Cauchy 动量残差显式约束动量平衡。然而，这一物理约束存在内在局限：依赖事先指定的本构律（流体/弹性/刚体）而非从数据中发现，对多材料耦合界面的应力连续性建模能力有限；Cauchy 约束本质上是"惩罚偏离"而非保证守恒结构，无法严格保证能量守恒；在光流监督失效的场景（遮挡、大运动、快速形变）下，物理约束缺乏独立的观测锚定。Baseline 问题。这种纯数据驱动的做法存在两个根本缺陷：(1) 将运动简化为刚体变换，忽略了不同材料（流体、弹性体、刚体）各自的运动本构律（constitutive law，即描述材料应力-应变关系的物理定律）；(2) 仅靠二维视觉投影无法统一约束三维粒子的运动状态，导致高斯粒子偏离物理上合理的运动轨迹——出现穿越、非物理拉伸、质量不守恒等伪影。核心思路。PIDG 的关键洞察是：每个高斯基元天然可以被视为一个拉格朗日物质粒子——就像在流体力学中跟踪流体微团的运动轨迹一样。因此，PIDG 将每个高斯的状态向量扩展为 $G_i = (\mu_i, q_i, s_i, c_i, o_i, e_i, v_i, \sigma_i)$，其中前三项是常规的位置、旋转（四元数）、缩放，$c_i, o_i$ 是颜色和不透明度，而 $e_i \in \mathbb{R}^H$（嵌入向量）、$v_i \in \mathbb{R}^{}3$（速度向量）和 $\sigma_i \in \mathbb{R}^{}6$（应力张量的六个独立分量）是新引入的独立物理属性。
\textbf{方法信息：} 代码：\url{https://github.com/scailab-ustc/physics-informed-deformable-gaussian-splatting}。
\textbf{设计动机与边界：} 相对原始 3DGS 以独立高斯和图像重建误差为主的设定，Physics-Informed Deformable Gaussian Splatting 把问题落在以下可实现的改动上：方法包含三个核心组件：(1)4D分解哈希编码：将4D时空坐标映射到四个3D哈希网格(Gxyz, Gxyt, Gyzt, Gxzt)，用空间MLP和注意力权重调制时间编码，结合两阶段优化策略分离静态和动态成分；(2)时变材料场：每个高斯被视为拉格朗日粒子，通过时变材料场独立预测速度和本构应力，将Cauchy动量方程的残差作为物理约束加入优化目标；(3)拉格朗日粒子流监督：将3D粒子运动投影到2D与光流匹配，加速收敛并提高泛化。其设计目的不是单独增加一个模块，而是用上述表示、监督或渲染约束消除原流程在目标任务中的具体失配。提出PIDG方法，将每个高斯粒子视为具有时变本构参数的拉格朗日物质点，通过Cauchy动量残差作为物理约束，独立预测每个粒子的速度和本构应力。采用静态-动态解耦的4D分解哈希编码，并利用拉格朗日粒子流与光流匹配进行监督，实现物理一致的单目动态重建。适用时还应保留论文给出的边界：物理约束基于简化的本构模型，对复杂非线性和多相材料建模能力有限；依赖光流估计质量；单目输入的尺度歧义问题。
% METHOD: 2605.11266
\paragraph{PG-3DGS（物理目标驱动的高斯优化）（2026.05）}
PG-3DGS~\cite{2605_11266} 是本节的另一条重要锚点：它不只是给重建后的高斯标注物理属性，而是在优化过程中让物理仿真反向影响高斯几何。普通 3DGS 从多视图图像中恢复外观时，可能得到外表看似合理、内部却封闭的茶壶，或外形像飞机但气动性能不成立的结构。PG-3DGS 的改变是把可微物理模拟接入 3DGS 的 inverse graphics 循环，同时优化视觉损失和物理目标。

具体做法是在每次迭代中先用当前高斯渲染多视图图像，计算与输入图像的视觉损失；再把高斯转换为可被模拟器使用的软占据 mask 或边界条件，送入可微流体或固体模拟，得到与任务相关的物理结果。物理目标可以是倒水任务中液体能否从壶嘴流出，也可以是飞行任务中物体能否产生足够升力。模拟器的梯度通过 Gaussian-to-mask 映射回传到高斯的均值、协方差和不透明度，驱动几何产生对物理有意义的改变。相对 PUGS 的属性预测，PG-3DGS 更强调“优化出来的形状必须能完成某个物理功能”；相对普通 3D 生成或重建，它把不可见的内部连通性、空腔、边界等图像损失不敏感的因素纳入训练。
\textbf{实现细节：} 沿"监督学习/前馈"路线，PIXIE 是首个通过监督学习训练通用神经网络预测 3D 场景物理属性的前馈框架，利用 CLIP 预训练特征把外观与物理特性关联；ReconPhys 是首个从单目视频联合推断物理属性和 3DGS 的前馈框架，推理速度小于 1 秒。沿"概率/不确定性"路线，PhysGS 把物理属性估计建模为 3DGS 上的贝叶斯推理问题，建模认知不确定性和偶然不确定性。仿真效率方面，FastPhysGS 提出实例感知粒子填充和双向图解耦优化，在 1 分钟内完成高保真物理仿真；PIG 通过 2D 像素到 3D 高斯的快速精确映射实现多材料耦合交互。
\textbf{技术实现补充：} 代表了一条更激进的路线，PG-3DGS（2026.05） 不再将物理仿真视为 3DGS 的后处理步骤，而是将可微分物理模拟直接嵌入 3DGS 的优化过程中，使生成的场景不仅视觉准确，还满足物理功能目标（如茶壶能倒水、机翼能产生升力）。与 PhysGaussian 的"先重建再仿真"范式不同，PG-3DGS 将物理仿真纳入重建目标函数本身。优化目标 $J(\theta; \mathcal{I}, \mathcal{P}) = L_{\mathrm{vis}}(\theta; \mathcal{I}) + L_{\mathrm{phys}}(\theta; \mathcal{P}) + L_{\mathrm{reg}}(\theta)$ 在视觉损失之外增加物理损失 $L_{\mathrm{phys}} = \mathcal{P}(\Phi(\theta))$，其中 $\Phi(\theta)$ 是从高斯几何到物理仿真网格的转换，$\mathcal{P}$ 是任务相关的物理功能评价（如倒出的液体体积、产生的升力）。直觉上，视觉损失只能约束外观——两个外观相同的茶壶，一个壶嘴连通内部可以倒水，另一个壶嘴只是外部凸起无法出水，仅靠图像无法区分。动机。现有 3DGS 重建方法仅优化视觉保真度，无法保证生成物体的物理功能性（如容器是否可盛水、机翼是否产生升力），因此需要将物理仿真约束嵌入重建目标函数中。PG-3DGS 首次在 3DGS 优化中引入这类物理约束，通过可微分流体仿真将物理梯度反传至高斯参数，使优化过程自动调整几何以满足功能需求。生成的场景可直接用于物理交互，无需人工后处理。紧耦合路线将高斯基元直接作为物理粒子，优势是无需格式转换，但限制为 MPM 一种物理引擎。许多实际场景需要更灵活的物理仿真器——例如刚体碰撞、有限元分析或商业游戏引擎。为此，松耦合路线通过引入代理结构（proxy structure）桥接高斯基元与任意物理引擎。
\textbf{方法信息：} 官方软件归档：\url{https://doi.org/10.5281/zenodo.19392335}。
\textbf{设计动机与边界：} 相对原始 3DGS 以独立高斯和图像重建误差为主的设定，PG-3DGS 把问题落在以下可实现的改动上：核心思想：标准3DGS仅优化视觉质量，忽略物理属性。PG-3DGS在3DGS优化循环中耦合可微分物理模拟器，同时优化视觉损失和物理目标。物理目标可包括稳定性（物体不悬浮）、碰撞响应（物体不穿透）等。梯度从物理模拟反向传播到高斯参数，使高斯在保持视觉质量的同时满足物理约束。其设计目的不是单独增加一个模块，而是用上述表示、监督或渲染约束消除原流程在目标任务中的具体失配。提出PG-3DGS，将可微分物理模拟与3DGS结合，使高斯表示不仅视觉准确，还满足物理目标（如稳定性、碰撞响应）。首次在3DGS优化中引入物理约束，使生成的场景可进行物理交互。适用时还应保留论文给出的边界：物理模拟的计算开销；对复杂物理现象（如流体）的支持有限。

\section{可变形物体与布料仿真}

可变形物体路线的关键矛盾是：3DGS 很适合渲染非刚性物体外观，但物理仿真需要拓扑、连接关系和材料参数。MaGS 用网格把高斯吸附到可变形结构上，PhysTwin 用弹簧质点系统构建可交互数字孪生，SketchRodGS 则把细长结构从高斯场中提取成可仿真的杆。
% METHOD: 2406.01593
\paragraph{MaGS（网格吸附高斯）（2024.06）}
MaGS~\cite{2406_01593} 是本路线的关键锚点。它针对的是 3DGS 与仿真结构之间的错位：纯高斯表示渲染灵活，却缺少网格拓扑；纯网格有拓扑和物理先验，却不如高斯容易拟合复杂外观。MaGS 的改变是构造 Mesh-adsorbed Gaussian，把每个高斯约束在三角网格附近游走，从而同时利用网格的可变形结构和高斯的渲染能力。

具体表示中，每个高斯除了普通均值、协方差、颜色和不透明度，还记录所在网格面片索引和局部逻辑坐标。空间位置通过面片顶点插值得到，Gaussian Hover 机制允许高斯沿法向有一定悬浮偏移，不必严格落在网格表面，这补偿了网格过粗或表面细节不足的问题。网络部分包含三个模块：RMD-Net 从视频中学习运动先验并精化网格变形；RGD-Net 建模网格与高斯之间的相对位移，提升外观拟合；MPE-Net 利用网格内在信息引导前两者，使方法能泛化到用户自定义变形，而不只复现训练视频。它可兼容 ARAP、SMPL 和软体物理仿真等先验。相对 GASP 的点式物理接口，MaGS 更强调“高斯要附着在有拓扑的可变形载体上”，因此适合人体、布料和软体对象。
\textbf{实现细节：} 前面讨论的 PIDG、NeHaD、PhysConvex、PhysFlow 等方法都在点云层面上引入物理先验——高斯基元在欧氏空间中自由运动，缺乏显式的曲面结构。另一条路线则尝试引入三角网格作为几何先验：DG-Mesh 和 SplattingAvatar 将高斯固定锚定在网格面片上，虽然获得了曲面信息，却牺牲了高斯本身的灵活性——高斯无法离开所在面片，也无法在面片间滑动；而 SC-GS 等纯点云方法在大幅编辑时则缺乏曲面连续性信息。MaGS（2024.06） （Mesh-adsorbed Gaussian Splatting）的核心设计思想介于两者之间：让高斯在网格面上"漫游"（roam）而非"锚定"（anchor），既保留网格的曲面连续性和模拟通用性，又不丧失高斯的灵活表达能力。Mesh-adsorbed Gaussian 表示。MaGS 的混合表示由三角网格和 3D 高斯两部分组成。三角网格具有 $N$ 个顶点和 $M$ 个面片，定义了顶点坐标 $v_i \in \mathbb{R}^{}3$、面片拓扑（每面片三个顶点索引）和法向量 $n_j \in \mathbb{R}^{}3$（可从顶点和面片计算得出，面片内任意位置的法向量通过三个顶点法向的线性插值获得（2024.06） ）。每个高斯在标准 3DGS 属性之外，额外存储两个关键属性：逻辑坐标 $w \in [0,1]^{}3$ 和所在面片编号 $\phi$。初始化时，$w$ 随机生成，$\phi$ 均匀分配到各面片以确保每个面片上高斯数量相等。高斯的空间位置通过面片顶点插值加法向偏移计算： 其中 $f$ 是重心插值函数 $f(w, x, y, z) = w_1 x + w_2 y + (1 - w_1 - w_2)z$，$l_j = |v_{j_1}v_{j_2} \times v_{j_1}v_{j_3}| \, / \, |v_{j_1}v_{j_2} \cdot v_{j_1}v_{j_3}|$ 是面片 $j$ 的缩放因子（反映面片面积的度量）。
\textbf{设计动机与边界：} 相对原始 3DGS 以独立高斯和图像重建误差为主的设定，MaGS 把问题落在以下可实现的改动上：MaGS的Mesh-adsorbed Gaussian将3D高斯约束在三角网格附近游走：每个高斯记录逻辑坐标w和所在面片索引phi，通过插值计算空间坐标。Gaussian Hover设计允许高斯不严格在网格表面而是悬浮于其上（通过法向偏移w\_3），增加了自由度。三个网络协同工作：RMD-Net从视频数据学习运动先验来精化网格变形；RGD-Net建模网格与高斯之间的相对位移以增强渲染保真度；MPE-Net利用网格内在信息引导前两个网络，使系统泛化到用户自定义变形（无需时序数据）。该框架兼容多种变形先验（ARAP、SMPL、软体物理仿真），在重建和仿真两个任务上同时取得SOTA性能。其设计目的不是单独增加一个模块，而是用上述表示、监督或渲染约束消除原流程在目标任务中的具体失配。提出Mesh-adsorbed Gaussian表示，首次在同一框架中兼顾3D重建的灵活性和仿真的结构性需求。通过Gaussian Hover设计增加高斯在网格上的自由度，结合三个专用网络（RMD-Net、RGD-Net、MPE-Net）实现从视频数据学习运动先序并进行用户自定义变形仿真，在D-NeRF、DG-Mesh和PeopleSnapshot数据集上达到SOTA。适用时还应保留论文给出的边界：依赖初始网格质量，对极端拓扑变化的处理有挑战。
\textbf{方法信息：} 代码：\url{https://github.com/wcwac/mags}；GitHub 关注度较高，约 104 stars。



% METHOD: 2503.17973
\paragraph{PhysTwin（可变形物体物理数字孪生）（2025.03）}
PhysTwin~\cite{2503_17973} 的锚点是从动态视频重建可变形物体。普通动态 3DGS 能复现观测运动，但一旦给物体新的推、拉、挤压动作，它没有内部物理模型来预测未来状态。PhysTwin 的改变是把弹簧-质点物理系统、生成式形状先验和 Gaussian Splatting 外观渲染组合成一个可交互数字孪生。

具体做法以多视角 RGBD 视频为输入，目标是同时估计几何、拓扑、物理参数和外观。第一阶段优化几何与物理：用生成式形状模型和视觉基础模型估计完整几何和拓扑，再把物体状态放入 spring-mass system，用显式 Euler 积分推进时间，并最小化预测 RGBD 与观测之间的差异。第二阶段固定或细化物理结构，优化高斯外观参数，使渲染与视频一致。由于状态演化来自显式物理模型，重建后的对象可以在新交互下实时模拟，例如绳索、毛绒玩具、布料和包裹的未来形变。相对 MaGS，PhysTwin 的重点更偏“从视频反推一个能仿真的物理系统”；相对普通物理引擎，它用高斯渲染把真实外观接到仿真状态上。
\textbf{实现细节：} PhysTwin（2025.03） 修改了可变形物体的动力学表示环节——用弹簧-质点图结构替代标准 3DGS 的无拓扑高斯基元集合，为离散高斯点注入结构化的物理连接关系。标准 3DGS 的动态变体（如 Deformable 3DGS、Dynamic 3DGS）通过学习形变场或逐帧优化高斯位置来拟合运动，但缺少物理约束使其无法外推到未见交互场景（如改变抓取方式或施加新的外力）。PhysTwin 的设计动机正是针对这一问题：以弹簧-质点模型作为物理骨架，每个质点通过弹簧与相邻质点连接，弹簧的弹力提供恢复力，阻尼力耗散能量，从而在拓扑层面建立可变形物体的内部结构。PhysTwin 的核心表示是一个弹簧-质点图 $G=(V,E)$，其中 $V$ 为质点集、$E$ 为弹簧集。质点 $i$ 受到相邻弹簧力和阻尼力的合力驱动： 其中 $k_{ij}$ 为弹簧刚度，$l_{ij}$ 为弹簧自然长度，$\gamma$ 为阻尼系数，$\mathbf{F}_i^{\mathrm{ext}}$ 包含重力、碰撞等外力。系统通过显式欧拉积分更新状态 $\mathbf{X}_{t+1} = f_{\boldsymbol{\theta}, G_0}(\mathbf{X}_t, \mathbf{a}_t)$，其中 $\boldsymbol{\theta}$ 编码所有物理参数（刚度、阻尼、碰撞参数），$G_0$ 为初始几何与拓扑，$\mathbf{a}_t$ 为控制点动作。直观而言，弹簧力将质点拉回平衡位置（类比橡皮筋），阻尼力防止永久振荡（类比减震器），二者共同约束物体的弹性变形行为。PhysTwin 的核心流程分两个阶段。第一阶段聚焦物理与几何优化：先用生成式形状先验（TRELLIS）从单张 RGB 图像恢复完整几何——因为稀疏视角下深度观测仅覆盖部分表面，直接从点云恢复完整拓扑不可行——获得初始网格后，通过配准模块（基于 SuperGlue 的 2D 匹配估计旋转、PnP 精修位姿、as-rigid-as-possible 非刚性形变）对齐到真实观测。
\textbf{实现补充：} 实现上还要核对输入、输出、参数、特征、初始化、训练损失、正则、渲染或推理接口，以及优化中的剪枝、分裂、采样、编码或解码步骤；这些抓手决定该方法到底改了 3DGS 管线的哪一环。
\textbf{锚点细化：} PhysTwin：针对 3DGS 高斯基元缺乏拓扑连接关系、动态变体通过学习形变场或逐帧优化高斯位置以拟合运动但缺少物理约束、无法外推到未见交互场景（如改变抓取方式或施加新外力）的局限，用\textbf{弹簧-质点图结构}替代无拓扑高斯基元集合，核心思想是\textbf{为离散高斯点注入结构化的物理连接关系}——弹簧力将质点拉回平衡位置（类比橡皮筋），阻尼力防止永久振荡（类比减震器），通过 TRELLIS 生成式形状先验从单张 RGB 图恢复完整几何，零阶 CMA-ES 采样优化不可微拓扑参数、一阶梯度精修稠密刚度，LBS 线性混合蒙皮将高斯绑定到质点网格，支持实时交互仿真和基于模型的机器人运动规划。
\textbf{方法信息：} 代码：\url{https://github.com/jianghanxiao/phystwin}；GitHub 关注度高，约 425 stars；引用很多，Semantic Scholar 约 85 次。



% METHOD: 2601.02072
\paragraph{SketchRodGS（从草图提取细长结构并动画化）（2026.01）}
SketchRodGS~\cite{2601_02072} 处理的是高斯场中常被忽视的细长几何，例如绳索、线缆、布料边缘。普通 3DGS 可以渲染这些结构，但没有显式中心线或杆模型，无法直接接入 elastic rod simulation。SketchRodGS 的改变是用用户草图作为交互信号，从 3DGS 中提取一条可仿真的折线骨架。

具体流程先修改高斯 rasterization，为每个像素记录贡献最大的高斯原语索引，形成 primitive index image。用户在屏幕空间画一笔标出细长部件，系统取得起点、终点和笔触附近像素对应的高斯。随后在屏幕空间搜索连接起终点的路径，路径代价同时考虑相邻高斯中心的三维距离和路径偏离用户笔触的二维距离；动态规划用于高效求解最短路径，并加入部分遮挡处理，避免路径被可见性中断。得到的折线可作为 elastic rod 或类似细长体的仿真骨架，再驱动高斯场中的相关原语运动。它相对 MaGS/PhysTwin 更轻量，不试图重建完整物理数字孪生，而是给已有 3DGS 场景补一个可编辑的细长结构抽取工具。
\textbf{实现细节：} SketchRodGS（2026.01） 从用户交互的角度解决 3DGS 场景中细长几何的提取与动画化问题。其核心思路十分直觉：用户在屏幕上画一条线（草图），系统自动找到对应的 3D 细长结构并将其转化为可仿真的折线网格（polyline mesh）。由于 3DGS 缺乏拓扑连接信息且高斯基元的配置包含大量噪声，直接从点云提取细长结构极为困难。SketchRodGS 的关键创新在于屏幕空间最短路径分析：将用户草图投影到多视角的深度图上，在屏幕空间中通过动态规划高效求解最短路径，从而在不同视角间鲁棒地重建出绳索、布料边缘等细长部件的折线表示。提取的折线可直接接入弹性杆（elastic rod）物理仿真器，使 3DGS 场景中的细长物体能够进行物理驱动的动画——绳索可以摆动、打结，布料边缘可以随风吹动。该方法在多个自然场景（in-the-wild）示例上验证了有效性。前述方法主要处理单时间步或短序列的物理行为。将物理仿真扩展到完整的时空 4D 场景，则需要解决材料参数的自动推断和跨帧运动一致性两个更根本的问题。下面介绍该方向的代表性方法。
\textbf{实现补充：} 实现上还要核对输入、输出、参数、特征、初始化、训练损失、正则、渲染或推理接口，以及优化中的剪枝、分裂、采样、编码或解码步骤；这些抓手决定该方法到底改了 3DGS 管线的哪一环。
\textbf{方法信息：} 代码：\url{https://github.com/haato-w/sketch-rod-gs}。
\section{流体与粒子动力学}

流体和粒子路线强调的是功能性重建：一个模型不仅要看起来像物体，还要在水流、烟雾、火焰等仿真中表现合理。FluidGaussian 把流体仿真不确定性用于主动重建，FieryGS 则把燃烧模拟与 3DGS 场景合成结合。


% METHOD: 2603.21356
\paragraph{FluidGaussian（基于流体仿真不确定性的主动重建）（2026.03）}
FluidGaussian~\cite{2603_21356} 的锚点是 ActiveNeRF/FisherRF 这类 next-best-view 规划。传统主动重建选择相机视角时主要看视觉不确定性，但一个几何在视觉上高分，仍可能因为孔洞、自交或薄结构错误而导致流体仿真失败。FluidGaussian 的改变是把流体-结构仿真的诊断结果作为视角选择标准，优先采集能同时提升视觉和物理保真度的视角。

具体流程是一个 plug-in。给定少量初始视图，先用视觉 NBV 方法提出 Top-K 候选相机姿态；然后对当前 3DGS 几何做流体冲刷或 SPH 风格的近表面仿真，计算仿真质量指标，例如近表面粒子速度散度。几何缺陷会在仿真中表现为异常流场、粒子穿透或不稳定区域；系统把这些仿真诱导的不确定性投回高斯或候选视图，对 Top-K 候选重新排序，选择最能暴露并修正功能缺陷的下一视角。相对 PG-3DGS，它不直接用物理梯度优化形状；相对普通 NBV，它让相机选择服务于“重建后能否做流体仿真”这一功能目标。
\textbf{实现细节：} 与 Gaussian Fluids 和 LagrangianSplats 均以仿真为核心目标不同，FluidGaussian（2026.03） 转向重建视角，提出了一个全新的思路：将流体仿真的不确定性引入 3DGS 重建。传统重建方法追求视觉质量，但重建结果在物理上可能不可行（如流体内部存在非零散度）。FluidGaussian 定义基于流体仿真的不确定性度量，将其与主动学习结合，优先选择同时提升视觉和物理保真度的视角。
\textbf{实现补充：} 实现上还要核对输入、输出、参数、特征、初始化、训练损失、正则、渲染或推理接口，以及优化中的剪枝、分裂、采样、编码或解码步骤；这些抓手决定该方法到底改了 3DGS 管线的哪一环。
\textbf{方法信息：} 代码：\url{https://github.com/delta-lab-ai/fluidgaussian}。
\textbf{设计动机与边界：} 相对原始 3DGS 以独立高斯和图像重建误差为主的设定，FluidGaussian 把问题落在以下可实现的改动上：FluidGaussian是一个即插即用方法，紧密耦合几何重建与流体-结构交互。核心步骤：(1)视觉不确定性提出Top-K候选相机姿态，采用ActiveNeRF或FisherRF作为基线NBV规划器；(2)FluidGaussian选择最佳物理信息相机姿态：使用SPH（光滑粒子流体动力学）对重建物体进行流体冲洗仿真，评估近表面流体粒子的仿真质量。不稳定或退化的近表面流体仿真暴露几何缺陷（间隙、自交、薄遮挡错误等）；(3)基于仿真诱导的不确定性度量从Top-K候选中选择最佳姿态，平衡视觉和物理质量。整个过程在有限视角预算下迭代进行。其设计目的不是单独增加一个模块，而是用上述表示、监督或渲染约束消除原流程在目标任务中的具体失配。提出FluidGaussian，首个将物理仿真不确定性引入3DGS重建的即插即用方法。定义了基于流体仿真的不确定性度量，并将其与主动学习结合。在NeRF Synthetic、Mip-NeRF 360和DrivAerNet++上实现PSNR提升8.6%、流体仿真速度散度降低62.3%的效果，开创了功能智能3D重建方向。适用时还应保留论文给出的边界：流体仿真增加了额外的计算开销；主要验证于物体级重建而非大尺度场景；流体-结构交互模型的选择可能影响通用性。
% METHOD: 2605.00177
\paragraph{FieryGS（物理集成火焰合成）（2026.05）}
FieryGS~\cite{2605_00177} 把 3DGS 用在真实场景火焰合成。普通视频生成或 2D 编辑能生成火焰外观，但很难保证火焰沿可燃材料传播、烟雾体积与场景几何对齐、碳化痕迹留在正确表面。FieryGS 的改变是把材料推理、体积燃烧仿真和 3DGS 渲染放在统一管线中。

具体做法先用 PGSR 等 3DGS 重建方法得到真实场景几何和外观，再对高斯进行聚类，把空间相近且视觉相似的高斯分成材料区域。每个区域渲染成 2D 视图后交给多模态大模型推断燃烧相关属性，例如材料类型、可燃性、热扩散和烟雾颜色。随后，Taichi 体积燃烧模拟在三维网格中计算火焰传播、热传导、烟雾生成和碳化。渲染时，系统把普通 3DGS rasterization 与参与介质体渲染结合，让火焰、烟雾和原场景同时显示。相对 FluidGaussian 关注重建质量如何影响流体仿真，FieryGS 更像把物理模拟作为可控视觉效果层接到高斯场上，用于 AR/VR、训练和安全演示。
\textbf{实现细节：} 与 FlameGS 聚焦于火焰光场的离线重建（从多视角投影中反推火焰三维光场分布）不同，FieryGS（2026.05） 从在线燃烧模拟角度出发，修改了 3DGS 场景编辑中的物理属性推理与特效生成方式——用 MLLM 驱动的零样本材料推理替代手工标注的可燃性参数，用 Taichi 体积燃烧模拟引擎替代预设视觉特效资产，再通过统一体积渲染器将模拟结果与 3DGS 场景融合。具体而言，框架先将重建的 3DGS 场景分割为语义一致区域，每个区域渲染为 2D 图像后送入 MLLM 推理材料类型、可燃性、热扩散系数和烟雾颜色等燃烧属性；随后 Taichi 引擎基于这些属性执行物理燃烧模拟，处理火焰沿可燃表面传播、烟雾体积扩散和表面碳化等复杂现象；最终统一体积渲染器将燃烧体素与 3DGS 高斯基元在同一管线中渲染，确保火光照明与场景几何的一致性。
\textbf{技术实现补充：} 与 FlameGS 聚焦于火焰光场的离线重建（从多视角投影中反推火焰三维光场分布）不同，FieryGS（2026.05） 从在线燃烧模拟角度出发，修改了 3DGS 场景编辑中的物理属性推理与特效生成方式——用 MLLM 驱动的零样本材料推理替代手工标注的可燃性参数，用 Taichi 体积燃烧模拟引擎替代预设视觉特效资产，再通过统一体积渲染器将模拟结果与 3DGS 场景融合。具体而言，框架先将重建的 3DGS 场景分割为语义一致区域，每个区域渲染为 2D 图像后送入 MLLM 推理材料类型、可燃性、热扩散系数和烟雾颜色等燃烧属性；随后 Taichi 引擎基于这些属性执行物理燃烧模拟，处理火焰沿可燃表面传播、烟雾体积扩散和表面碳化等复杂现象；最终统一体积渲染器将燃烧体素与 3DGS 高斯基元在同一管线中渲染，确保火光照明与场景几何的一致性。直觉上，可以理解为三个模块各司其职：MLLM 充当"材料专家"，自动判断场景中每个物体能否燃烧、如何燃烧；物理引擎充当"燃烧模拟器"，负责火焰传播和烟雾扩散的物理计算；渲染器充当"画面合成器"，负责将模拟结果无缝融合到 3DGS 场景中。这是将物理准确燃烧模拟集成到 3DGS 的首个工作，支持火焰强度、气流和点火位置等参数的用户交互控制。
\textbf{方法信息：} 项目页：\url{https://pku-vcl-geometry.github.io/FieryGS/}；当前未确认完整公开代码。
\textbf{设计动机与边界：} 相对原始 3DGS 以独立高斯和图像重建误差为主的设定，FieryGS 把问题落在以下可实现的改动上：FieryGS物理驱动3DGS火灾合成：(1)MLLM推理材料燃烧属性+HDBSCAN聚类分割可燃/不可燃高斯组；(2)Taichi 256\^{}3体积燃烧模拟——火焰传播、热传递、烟雾生成；(3)统一渲染器融合3DGS光栅化与体积燃烧渲染。其设计目的不是单独增加一个模块，而是用上述表示、监督或渲染约束消除原流程在目标任务中的具体失配。将物理准确燃烧模拟集成到3DGS，MLLM自动推理材料属性，支持火焰传播/烟雾扩散/碳化精细控制。适用时还应保留论文给出的边界：256\^{}3模拟网格分辨率限制细粒度火焰细节；MLLM材料推理可能对罕见材料推断不准确。

\section{其他物理方法}

这一节包含更专门的物理现象建模。与前面偏“动态视觉拟合”的方法不同，这里关心高斯能否承载真实物理状态：风场、动量交换、材料反应、火焰烟雾或连续场变量。DiffWind 说明 3DGS 不只可以表现物体表面，也可以作为风-物体交互中的拉格朗日粒子表示，与欧拉风场和流体约束共同优化；这类方法的重点是让可渲染高斯进入可控仿真，而不是只复现训练视频。
% METHOD: 2603.09668
\paragraph{DiffWind（风驱动物体的可微物理建模）（2026.03）}
DiffWind~\cite{2603_09668} 的锚点是动态 3DGS 对风吹物体的外观拟合。普通动态高斯可以学到树叶、布料或轻物体在视频中的运动，但它不显式知道风场，也不能重定向风力。DiffWind 的改变是把物体表示成 3DGS 粒子系统，把风表示成欧拉网格物理场，并通过 MPM 在粒子和网格之间交换动量。

具体管线先重建物体的高斯表示，并加入各向异性约束、不透明度约束、不可见区域 densification 和体积保持损失，减少针状伪影、漂浮点和弱纹理空洞。风场则在背景网格节点上存储密度、速度和力。MPM 的 P2G 步骤把高斯粒子的质量和速度传到网格，网格上施加风力并更新流场；G2P 步骤再把速度和位置更新传回粒子。训练时，系统通过可微渲染和可微模拟联合优化风场与物体运动，并加入 Lattice Boltzmann Method 约束，让风场符合基本流体动力学规律。这样得到的模型不仅能重现观测视频，还能做风力重定向：改变风的方向或强度，预测物体新的动态响应。
\textbf{实现细节：} 与前述方法均针对单一物理现象（火焰、降雨）不同，DiffWind（2026.03） 聚焦于风-物体交互这一跨模态耦合问题，修改了动态场景重建中的运动建模方式——用可微分物理模拟驱动的风场-物体联合优化替代纯数据驱动的变形场学习，提出首个物理信息可微分风-物体交互建模框架。其核心创新在于为风和物体分别选择最适合的表示并进行耦合交互：风被建模为规则的网格物理场，每个网格节点存储密度、速度和力等物理量，用 Lattice Boltzmann Method（LBM）模拟其演化以符合流体力学基本定律；物体则继承 3DGS 的粒子表示，每个高斯基元携带外观、材质和运动属性；两者之间的交互通过 Material Point Method（MPM）实现——MPM 将粒子属性映射到背景网格上计算受力，再将更新后的物理量传回粒子，自然桥接了连续场（风）与离散粒子（物体）两种表示。在重建阶段，通过可微分模拟与可微分渲染的端到端反向传播，联合优化时空风场和物体运动参数；在 LBM 物理约束的保障下，优化得到的风场不仅匹配观测运动，还满足流体守恒律。该框架支持前向模拟（在新风条件下预测物体运动）和风力重定向（将重建的风场效果迁移到新场景），并构建了 WD-Objects 风驱动场景数据集用于系统评测。
\textbf{实现补充：} 实现上还要核对输入、输出、参数、特征、初始化、训练损失、正则、渲染或推理接口，以及优化中的剪枝、分裂、采样、编码或解码步骤；这些抓手决定该方法到底改了 3DGS 管线的哪一环。
\textbf{方法信息：} Ant Group / Alibaba 参与了这篇论文；代码：\url{https://github.com/zju3dv/DiffWind}。
\begin{takeaway}
本章的核心演进是从“高斯能看”走向“高斯能用”。PUGS 让高斯携带质量、摩擦等物理属性，PG-3DGS 让可微物理目标反向塑造高斯几何，PIDG 把动态高斯放进本构约束，GASP 提供与黑盒物理引擎连接的实用接口。可变形物体方向中，MaGS 用网格吸附补拓扑，PhysTwin 用弹簧质点构建数字孪生，SketchRodGS 用草图抽取可仿真的细长杆。流体与风火方向则进一步说明，物理感知 3DGS 不只是让动态画面更真实，而是让重建结果能参与流体、燃烧、风场和功能性仿真。
\end{takeaway}
